% Source: https://artofproblemsolving.com/community/c3h2222832_100th_post_celebration_avatar_themed
% Stable identity: aops-topic-2222832-q01
Seen from above, the tart is the region between circles of radius $\dfrac{3}{2}$ and $\dfrac{1}{2}$ (outer diameter $3$, width $1$). Its cross-sections are discs of radius $\dfrac{1}{2}$ whose centres lie on the circle of radius $R=\dfrac{3}{2}-\dfrac{1}{2}=1$ (dashed).

\begin{asy}
usepackage("amsmath");
size(5cm);
defaultpen(fontsize(10pt));
pair O=(0,0);
filldraw(circle(O,1.5),lightgray,black);
filldraw(circle(O,0.5),white,black);
draw(circle(O,1),dashed+blue);
draw((-1.5,0)--(1.5,0),Arrows(5));
label("$3$",(0.9,0),N);
draw((0.5,-0.4)--(1.5,-0.4),red,Arrows(4)); label("$1$",(1,-0.4),S,red);
draw(O--1*dir(130),blue,Arrow(4)); label("$1$",0.5*dir(130),SW,blue);
dot(O);
\end{asy}

The torus is swept out by rotating a disc of radius $\rho=\dfrac{1}{2}$ about an axis at distance $R=1$ from the disc's centre. By Pappus' theorem the volume is the area of the disc times the distance travelled by its centre:
\[V=\pi\rho^2\cdot2\pi R=\pi\cdot\dfrac{1}{4}\cdot2\pi\cdot1=\boxed{\dfrac{\pi^2}{2}}\ \text{cubic inches}\approx4.93\text{ cubic inches}.\]
